% generated from Bio 1/dbg_sandpile.tex
\chapter{Assembly ambiguity is not determined by branch statistics: critical groups of de Bruijn graphs}\label{chap:Bio1}
\chaptermark{Bio1}
\begin{chapterabstract}
\noindent
The number of arc-rooted circular sequences sharing a given $k$-mer spectrum is, by the
BEST theorem, the product of a \emph{local} factor $\prod_v(d^+(v)-1)!$, which depends only
on the out-degree multiset of the de Bruijn graph, and a \emph{global} factor, which is the
order of the graph's critical (sandpile) group. We make four points about the global factor.
First, it is not determined by the degree data: we exhibit two sequences of length $60$
whose de Bruijn graphs at $k=4$ have identical vertex counts, identical out-degree
multisets and hence identical local factors, but admit $4$ and $256$ reconstructions --- a
factor of $64$. This is a strict sharpening of the observation of Kingsford, Schatz and Pop
that repeat counting is insufficient. Second, the group structure reads off repeat
parameters: a word of length $L$ occurring $r$ times in an otherwise $k$-unique sequence
contributes $\K(G_k)\cong(\Z/r)^{L-k}$. Third --- and this is new in this version --- the
group detects how repeats \emph{interleave}: for two repeats of length $L$ occurring twice
each, the disjoint and nested arrangements $A\,A\,B\,B$ and $A\,B\,B\,A$ have isomorphic
compacted graphs and $\K(G_k)\cong(\Z/2)^{2s-2}$, the direct sum of the two single-repeat
contributions, while the interleaved arrangement $A\,B\,A\,B$ has $\K(G_k)\cong(\Z/2)^{2s-1}$,
one extra $\Z/2$ that survives to $k=L$ after both individual contributions have died. The
degree data are identical in all three cases. Fourth, the invariant is cheap where it
matters: contracting unitigs leaves the critical group unchanged, so it is computed on the
compacted graph. We demonstrate this on a real genome: for bacteriophage $\varphi$X174
(5386 bp) the $k=10$ graph has $5339$ vertices but $47$ branch points, and its critical
group, of order $4.4\times10^{12}$, is computed in ten milliseconds; the spectrum collapses
through $132$, $2$, $1$ at $k=11,12,13$, one past the longest exact repeat. All claims are
verified by machine: $30$ checks, exact integer arithmetic except where labelled.
\end{chapterabstract}

\section{Introduction}\label{Bio1:sec:intro}

Given a collection of $k$-mers with multiplicities --- an idealised sequencing experiment
--- how many genomes are consistent with it? The $k$-mers and their overlaps form a de
Bruijn multigraph $G_k$ in which the genome is an Eulerian circuit, so the answer is the
number of Eulerian circuits, which the BEST theorem \cite{AEdB,SmithTutte} evaluates as
\begin{equation}\label{Bio1:eq:best}
\ec(G_k)=\bigl|\K(G_k)\bigr|\cdot\prod_{v}\bigl(d^{+}(v)-1\bigr)!\,,
\end{equation}
$\K(G_k)$ being the group of spanning arborescences --- the \emph{critical} or
\emph{sandpile} group. Kingsford, Schatz and Pop \cite{KSP} computed this for real
prokaryotic genomes as a measure of assembly complexity; Loukides, Pissis and co-authors
\cite{Loukides} gave fast algorithms for assessing it.

Equation~\\eqref{Bio1:eq:best} splits the ambiguity into two factors of different character. The
product of factorials is a function of the out-degree multiset alone --- the \emph{branch
statistics}, read off the graph by inspection. The other factor is not. This paper is
about what the other factor knows.

\paragraph{What is already known.}
Two literatures meet here, and we claim nothing in either. On the combinatorial side the
critical groups of \emph{complete} de Bruijn graphs are completely understood: Knuth
\cite{Knuth} gave the line-digraph recursion, Levine \cite{Levine} the group structure
$\K(\DB(2,n))\cong\bigoplus_{j=1}^{n-1}(\Z/2^{j})^{2^{\,n-1-j}}$, and Chan, Hollmann and
Pasechnik \cite{CHP} the generalised de Bruijn and Kautz cases. On the assembly side,
counting reconstructions via BEST is established \cite{KSP,Loukides}, as are uniqueness and
safety criteria \cite{ObscuraTomescu} and feasibility thresholds \cite{BBT,NagarajanPop}.
What appears not to have happened is that the two ever met: we have found no work applying
the critical group of a graph to assembly, and none that looks at the \emph{structure} of
$\K(G_k)$ for a genome de Bruijn graph, which is a subgraph of a complete one and not
covered by \cite{Levine,CHP}.

\paragraph{What is new here.}
\begin{enumerate}
\item[(i)] \textbf{Theorem~\\ref{Bio1:thm:notdeg}}: the global factor is not a function of the
degree data. Two explicit sequences have identical branch statistics and reconstruction
counts differing by $64\times$.
\item[(ii)] \textbf{Theorem~\\ref{Bio1:thm:repeat}}: for a single repeat of length $L$ and
multiplicity $r$, $\K(G_k)\cong(\Z/r)^{L-k}$.
\item[(iii)] \textbf{Theorem~\\ref{Bio1:thm:interleave}} (new in this version): for two repeats
the group is additive in the disjoint and nested arrangements and picks up exactly one
extra $\Z/2$ in the interleaved arrangement. Since interleaved repeats are the notoriously
hard case for assembly, this is the first place the algebra and the practice visibly
agree. Proposition~\\ref{Bio1:prop:three} then shows that for three repeats the correction is
the sandpile group of a small digraph and not a chord-diagram rank --- it can be $\Z/3$.
\item[(iv)] \textbf{Lemma~\\ref{Bio1:lem:contract}} makes the invariant computable on the
compacted graph, and Section~\\ref{Bio1:sec:real} computes it on a real genome.
\end{enumerate}

\paragraph{Relation to \cite{KSP}.}
Kingsford, Schatz and Pop already argue that repeat counting is insufficient, with the
example that three repeats in two copies each ($8$ reconstructions) is easier than one
repeat in six copies ($720$). Their two situations have \emph{different} out-degree
multisets, $(2,2,2)$ against $(6)$, so the local factor already distinguishes them.
Theorem~\\ref{Bio1:thm:notdeg} fixes the out-degree multiset and still finds a $64$-fold spread,
and Theorem~\\ref{Bio1:thm:interleave} says what the residual information is: not how many
repeats or how many copies, but how they are wired to one another.

\section{Setting}\label{Bio1:sec:setting}

Let $S$ be a circular sequence of length $N$ over an alphabet of size $\sigma$. For $k\ge1$
the \emph{de Bruijn multigraph} $G_k=G_k(S)$ has as vertices the distinct $k$-mers of $S$
and one arc for each of the $N$ occurrences of a $(k{+}1)$-mer, from its length-$k$ prefix
to its length-$k$ suffix. $S$ traverses every arc once, so $G_k$ is connected and balanced,
and $S$ is an Eulerian circuit; every Eulerian circuit is a circular sequence with the
same $k$-mer spectrum, counted with multiplicities. We write $n_k=|V(G_k)|$, $d^{+}(v)$ for
out-degree (the number of occurrences of $v$), and $\Loc(G):=\prod_v(d^{+}(v)-1)!$.

\begin{remark}[Counting convention]\label{Bio1:rem:convention}
Throughout, reconstruction counts are \emph{arc-rooted}: circuits are counted as beginning
with a fixed arc, which is the convention under which \\eqref{Bio1:eq:best} holds
\cite{AEdB,SmithTutte}. Circular sequences up to cyclic shift are counted by dividing by
$N$ when $S$ is aperiodic, and by Burnside's lemma otherwise; none of the algebra below
depends on which convention is chosen, since the correction is a function of $N$ alone.
\end{remark}

The \emph{critical group} is $\K(G):=\coker(\Delta^{\mathsf T})$, $\Delta$ the out-Laplacian
$D^{+}-A$ with the row and column of a fixed root deleted; for a connected balanced digraph
the isomorphism class does not depend on the root, and $|\K(G)|$ is the number of spanning
arborescences. We compute it as a Smith normal form.

\begin{remark}[The tower]\label{Bio1:rem:tower}
$G_{k+1}(S)$ is a spanning subgraph of the line digraph $L(G_k(S))$, with equality for the
complete de Bruijn graph; the prefix and suffix maps are the tail and head maps of $G_k$
and are the degeneracy maps of the non-normal tower of \cite[V]{ANCFT}. The line-digraph
recursion of \cite[III]{ANCFT} specialises to Knuth's formula and reproduces
$|\K(\DB(\sigma,k))|=\sigma^{\sigma^{k}-k-1}$ (check~(T)); check~(L) reproduces Levine's
decomposition. These are consistency checks on our conventions, not results.
\end{remark}

\section{The local and global factors}\label{Bio1:sec:decomp}

\begin{proposition}[Decomposition]\label{Bio1:prop:decomp}
The number of arc-rooted circular sequences with the same $k$-mer spectrum as $S$ equals
$|\K(G_k)|\cdot\Loc(G_k)$. The factor $\Loc(G_k)$ depends only on the multiset
$\{d^{+}(v)\}$; the factor $|\K(G_k)|$ does not.
\end{proposition}

The first sentence is \\eqref{Bio1:eq:best}. The second is Theorem~\\ref{Bio1:thm:notdeg}.

\section{Contraction, and a single repeat}\label{Bio1:sec:repeat}

\begin{lemma}[Unitig contraction]\label{Bio1:lem:contract}
Let $v$ be a vertex of a connected balanced multidigraph $G$ with $d^{-}(v)=d^{+}(v)=1$,
whose incoming and outgoing arcs are distinct, say $u\to v$ and $v\to w$. Let $G'$ be
obtained by deleting $v$ and replacing those two arcs by a single arc $u\to w$. Then
$\K(G')\cong\K(G)$. Self-loops may be deleted freely: a loop at $v$ adds $1$ to both
$D^{+}_{vv}$ and $A_{vv}$ and so does not appear in $\Delta$.
\end{lemma}

\begin{proof}
Root away from $v$ and $u$. The row of $v$ in $\Delta$ is $e_v-e_w$ and the column of $v$
has a single off-diagonal entry, in the row of $u$; adding the row of $v$ to the row of $u$
clears it and turns the row of $u$ into its row in $\Delta(G')$, after which the row and
column of $v$ form a unit block. Both operations are unimodular. The statement about loops
is immediate from $\Delta=D^{+}-A$ and is check~(P).
\end{proof}

Iterating the lemma yields the \emph{compacted} (unitig) de Bruijn graph, in which every
maximal non-branching path is one arc and only branch vertices survive.

\begin{theorem}[A repeat of length $L$ and multiplicity $r$]\label{Bio1:thm:repeat}
Let $W$ have length $L$ and let $S$ be circular with every $k$-mer occurring once except
the $L-k+1$ $k$-mers of $W$, each occurring exactly $r\ge2$ times. Then for $1\le k\le L$
\[
\K\bigl(G_k(S)\bigr)\;\cong\;(\Z/r)^{\,L-k}.
\]
\end{theorem}

\begin{proof}
Put $s=L-k+1$ and let $v_1,\dots,v_s$ be the $k$-mers of $W$ in order. Every other vertex
has in- and out-degree $1$; contract them. What remains is the spine $v_i\to v_{i+1}$ with
multiplicity $r$ and $r$ arcs $v_s\to v_1$. Rooted at $v_1$, on $v_2,\dots,v_s$ the
out-Laplacian is $r(I-N)$ with $N$ the nilpotent shift; $I-N$ is unimodular, so the
cokernel is $(\Z/r)^{s-1}$.
\end{proof}

\section{Two repeats: the group sees interleaving}\label{Bio1:sec:interleave}

Fix two distinct words $A,B$ of length $L$ and a circular sequence in which each occurs
exactly twice, non-overlapping, with every other $k$-mer unique; put $s=L-k+1$. The four
occurrences can be arranged in three ways around the circle:
\[
\text{disjoint}\ A\,A\,B\,B,\qquad \text{nested}\ A\,B\,B\,A,\qquad
\text{interleaved}\ A\,B\,A\,B.
\]
The $k$-mer spectra of the three sequences have identical multiplicity profiles, so
$n_k$, the out-degree multiset, and hence $\Loc(G_k)$ agree. This is Theorem~\\ref{Bio1:thm:notdeg}'s
setting one level up, and the question is whether $\K$ can tell them apart.

\begin{theorem}[Two repeats]\label{Bio1:thm:interleave}
With $a_1,\dots,a_s$ and $b_1,\dots,b_s$ the $k$-mers of $A$ and $B$:
\begin{enumerate}
\item[(i)] the disjoint and nested arrangements have isomorphic compacted graphs, and
\[
\K(G_k)\;\cong\;(\Z/2)^{2s-2}\;=\;(\Z/2)^{s-1}\oplus(\Z/2)^{s-1},
\]
the direct sum of the two single-repeat groups of Theorem~\\ref{Bio1:thm:repeat};
\item[(ii)] the interleaved arrangement has
\[
\K(G_k)\;\cong\;(\Z/2)^{2s-1}.
\]
\end{enumerate}
In particular the arrangement is detected by $\K$ and not by the degree data, and at
$k=L$ ($s=1$) the interleaved arrangement retains a $\Z/2$ after both single-repeat
contributions have become trivial.
\end{theorem}

\begin{proof}
Contract by Lemma~\\ref{Bio1:lem:contract}. The surviving vertices are the $2s$ shared $k$-mers,
each of in- and out-degree $2$. The spine arcs $a_i\to a_{i+1}$ and $b_i\to b_{i+1}$ have
multiplicity $2$. The four spacer paths contract to four connector arcs from
$\{a_s,b_s\}$ to $\{a_1,b_1\}$, one per transition between consecutive occurrences around
the circle. Reading the cyclic order: disjoint gives $a_s\!\to\!a_1$, $a_s\!\to\!b_1$,
$b_s\!\to\!b_1$, $b_s\!\to\!a_1$; nested gives $a_s\!\to\!b_1$, $b_s\!\to\!b_1$,
$b_s\!\to\!a_1$, $a_s\!\to\!a_1$ --- the same multiset, so the two compacted graphs
coincide; interleaved gives $a_s\!\to\!b_1$ twice and $b_s\!\to\!a_1$ twice.

(ii) Root at $a_1$. On $a_2,\dots,a_s,b_1,\dots,b_s$ in this order the reduced out-Laplacian
is $2(I-N)$, $N$ the shift along that path (both arcs of $b_s$ go to the deleted root).
$I-N$ is unimodular, so $\K\cong(\Z/2)^{2s-1}$.

(i) Root at $a_1$; write $x_v$ for the generator of the cokernel presentation at $v$. The
rows of $a_s$ and $b_s$ read $2x_{a_s}-x_{b_1}$ and $2x_{b_s}-x_{b_1}$, so $x_{b_1}=2x_{a_s}$
and $x_{b_1}$ can be eliminated. On the remaining $2s-2$ generators
$a_2,\dots,a_s,b_2,\dots,b_s$ every relation is now even --- the spine relations
$2x_{a_i}-2x_{a_{i+1}}$ and $2x_{b_i}-2x_{b_{i+1}}$, the relation $2x_{b_s}-2x_{a_s}$, and the
former row of $b_1$, which becomes $4x_{a_s}-2x_{b_2}$ --- so the relation matrix is $2M'$
with
\[
M'=\begin{pmatrix}
\text{rows } e_{a_i}-e_{a_{i+1}},\ 2\le i\le s-1\\
e_{b_s}-e_{a_s}\\
2e_{a_s}-e_{b_2}\\
\text{rows } e_{b_i}-e_{b_{i+1}},\ 2\le i\le s-1
\end{pmatrix}.
\]
Expanding $\det M'$ along the columns $a_2,\dots,a_{s-1}$ in turn (each has a single entry,
$+1$) leaves the $s\times s$ block on columns $(a_s,b_2,\dots,b_s)$ with rows
$2e_{a_s}-e_{b_2}$, $e_{b_i}-e_{b_{i+1}}$ ($2\le i\le s-1$), $e_{b_s}-e_{a_s}$.
Expanding along the first row with the cofactor signs written out,
$\det M'=2\,(-1)^{1+1}M_{11}+(-1)\,(-1)^{1+2}M_{12}=2M_{11}+M_{12}$. The minor $M_{11}$ (the
remaining rows on columns $b_2,\dots,b_s$) is upper bidiagonal with unit diagonal, so
$M_{11}=1$; the minor $M_{12}$ (the same rows on columns $a_s,b_3,\dots,b_s$) has a single
entry $-1$ in its first column, in its last row, and expanding along that column leaves a
lower bidiagonal block with diagonal $-1$, the signs combining to $M_{12}=-1$. Hence
$\det M'=2-1=1$ in this column order, $|\det M'|=1$ in any order, $M'$ is unimodular and
$\K\cong(\Z/2)^{2s-2}$.
\end{proof}

\begin{remark}[Reading]\label{Bio1:rem:reading}
For a single repeat the group is trivial at $k=L$: the repeated $L$-mer is a branch vertex
but its two exits lead to different places and there is only one way to thread it. For
two interleaved repeats at $k=L$ there is a genuine choice --- leave $A$ for the first $B$
or the second --- and $\Z/2$ records it. Nested and disjoint repeats offer no such
choice: whichever way one threads $A$, the two $B$'s are traversed consecutively. This is
the algebraic form of the folklore that interleaved repeats are what makes assembly hard.
Additivity of $\K$ over repeats, asked for in version~1, therefore holds exactly when the
repeats do not interleave; the general correction term is indexed by the interleaving
pattern, and Theorem~\\ref{Bio1:thm:interleave} is its two-repeat case.
\end{remark}

Check~(I) verifies the theorem three ways: on clean random sequence instances (the
prediction is exact in every one), on the explicit compacted graphs (six branch vertices;
connector multiplicities $\{1,1,1,1\}$ for disjoint and nested and $\{2,2\}$ for
interleaved), and by symbolic Smith normal form of the matrices in the proof for
$s\le7$.

\subsection*{Three repeats, and what the correction is not}

With three repeats the natural guess --- suggested to us by a reader --- is that the
interleaving correction is a function of the chord diagram of the six occurrences through
the rank over $\mathbb F_2$ of its intersection matrix. It is not.

\begin{proposition}[Three repeats]\label{Bio1:prop:three}
Let $A,B,C$ be distinct words of length $L$, each occurring exactly twice, with every other
$k$-mer unique and $s=L-k+1$. Let $D_w$ be the digraph on the three symbols with one arc
from the symbol of each occurrence to the symbol of the next occurrence around the circle,
loops deleted. Then $\K(G_L)\cong\K(D_w)$, and according to which pairs of the three chords
cross:
\begin{center}\small
\begin{tabular}{lccc}
\toprule
crossing pairs & $\mathbb F_2$-rank of intersection matrix & $\K(G_k)$, $s=1$ & $\K(G_k)$, $s=2$\\
\midrule
none            & $0$ & $0$            & $(\Z/2)^{3}$\\
one pair        & $2$ & $\Z/2$         & $(\Z/2)^{4}$\\
a path of two   & $2$ & $\Z/3$         & $(\Z/2)^{3}\oplus\Z/3$\\
all three pairs & $2$ & $(\Z/2)^{2}$   & $(\Z/2)^{5}$\\
\bottomrule
\end{tabular}
\end{center}
In every case $\K(G_k)\cong(\Z/2)^{3(s-1)}\oplus\K(D_w)$ for $s=1,2$.
\end{proposition}

\begin{proof}
At $s=1$ the compacted graph is $D_w$ by construction (Lemma~\\ref{Bio1:lem:contract} deletes
the loops), which gives $\K(G_L)=\K(D_w)$; the four values are the sandpile groups of the
four $2$-regular digraphs on three vertices that occur. For instance the path pattern
$A\,B\,A\,C\,B\,C$ gives the complete symmetric digraph on three vertices, whose sandpile
group is $\Z/3$. The $s=2$ column and the splitting are verified exactly on all $90$
circular words with two of each letter (check~(M)); we have not proved the splitting for
general $s$.
\end{proof}

\begin{remark}\label{Bio1:rem:notrank}
Three of the four patterns have intersection rank $2$ and give three different groups,
one of them of odd order: multiplicity-two repeats can interleave into $3$-torsion. So the
correction is not a chord-diagram rank, nor even a $2$-group. The right object is the
sandpile group of the connector digraph $D_w$, which for two repeats is trivial or $\Z/2$
(Theorem~\\ref{Bio1:thm:interleave}) and for three is the table above.
\end{remark}

\section{Branch statistics are insufficient}\label{Bio1:sec:notdeg}

\begin{theorem}[The global factor is not degree-determined]\label{Bio1:thm:notdeg}
There exist circular sequences $S_{\mathrm{lo}},S_{\mathrm{hi}}$ of length $60$ over
$\{A,C,G,T\}$ such that at $k=4$ their de Bruijn graphs have the same number of vertices
$n_4=51$ and the same out-degree multiset (seven branch vertices of out-degree $2$ and one
of out-degree $3$), hence the same local factor $\Loc=2$, but
\[
\K(G_4(S_{\mathrm{lo}}))=\Z/2,\qquad
\K(G_4(S_{\mathrm{hi}}))=(\Z/2)^{4}\oplus\Z/8 ,
\]
so that the numbers of reconstructions are $4$ and $256$. Explicit witnesses:
\begin{center}\ttfamily\small
\begin{tabular}{ll}
$S_{\mathrm{lo}}$: & GAACGAGAGAAAAACACGTACCAGCGCGGGGGGCGGGACGGCAGGTGACGCTCTGGGAAG\\
$S_{\mathrm{hi}}$: & GCACAGGGTCCAGTCTGCGAACGCGCTACAGTCTGGGTTCAAGCGGTCCTCTGTGCATCG\\
\end{tabular}
\end{center}
\end{theorem}

\begin{proof}
Exact computation; check~(N).
\end{proof}

\begin{corollary}\label{Bio1:cor:insufficient}
No function of the out-degree multiset of $G_k$ determines the number of reconstructions.
\end{corollary}

In a randomised search over $1500$ sequences of length $60$ at $k\in\{4,5\}$, $67$ of the
$98$ degree-classes containing more than one sequence contained more than one critical
group (check~(N), sampled).

\section{Computation}\label{Bio1:sec:compute}

Computing $\K(G_k)$ directly is one Smith normal form of an $(n_k-1)\times(n_k-1)$ integer
matrix, $O(n_k^{3})$ with the usual bit-length control \cite[I]{ANCFT}, which for a genome
is impractical. Lemma~\\ref{Bio1:lem:contract} removes the difficulty: the critical group of $G_k$
equals that of its compacted graph, whose vertices are the branch points. The compaction is
linear in $N$ --- one pass to find the branch vertices, one walk per outgoing arc --- and
assemblers already build this graph. The cost is therefore governed by repeat content, not
genome length; the next section shows what this means in practice.

\section{A real genome}\label{Bio1:sec:real}

Bacteriophage $\varphi$X174 (GenBank NC\_001422.1, $5386$ bp, circular) was the first
genome ever sequenced. Its longest exact circular repeat has length $R=12$.
Table~\\ref{Bio1:tab:phix} gives the critical-group spectrum around that scale. The compaction
is dramatic: at $k=10$ the graph has $5339$ vertices and $47$ branch points, and the Smith
normal form takes ten milliseconds. The group order at $k=10$, $4.4\times10^{12}$, is the
global part of the number of $10$-mer reconstructions of the phage; it falls to $132$ at
$k=11$, to $2$ at $k=12$, and the graph becomes a single cycle at $k=13=R+1$.

\begin{table}[ht]
\centering\small
\begin{tabular}{rrrrl}
\toprule
$k$ & $n_k$ & branch vertices & $|\K(G_k)|$ & invariant factors\\
\midrule
$\le9$ & $\le5233$ & $\ge146$ & --- & above the exact-SNF cap used here\\
10 & 5339 & 47 & $4\,408\,682\,819\,584$ & $2^{8},\dots,4305354316$\\
11 & 5376 & 10 & $132$ & $2,\ 66$\\
12 & 5384 &  2 & $2$   & $2$\\
13 & 5386 &  0 & $1$   & --- (single cycle)\\
\bottomrule
\end{tabular}
\caption{$\varphi$X174, NC\_001422.1. Branch vertices are the vertices of the compacted
graph. All entries exact (check~(F)); the sequence file is distributed with the script and
no network access is needed.}
\label{Bio1:tab:phix}
\end{table}

Two remarks. The levels $k\le9$ have compacted graphs of $146$ to $1384$ vertices; they are
within reach of a sparse or $p$-adic Smith normal form but not of the dense routine used
here, and we have not pursued them. And $k^{*}=R+1$ for this genome. A single isolated repeat collapses at $k=R$
(Theorem~\\ref{Bio1:thm:repeat}), so a surviving $\Z/2$ at $k=12$ is \emph{consistent with}
Theorem~\\ref{Bio1:thm:interleave} rather than implied by it; the companion paper \cite{Bio2}
settles the question by exhibiting the two reconstructions: the longest $12$-mer repeat is
interleaved with a second repeated $12$-mer, and the residual bit is exactly the one of
Theorem~\\ref{Bio1:thm:interleave} at $s=1$.

\section{What is open, and what is false}\label{Bio1:sec:open}

\paragraph{False.} The collapse order is not one more than the longest repeat. With $R(S)$
the longest repeated substring and $k^{*}(S)$ the least $k$ with $\K(G_k)$ trivial, over
$40$ random sequences of length $100$ we find $k^{*}=R$ in $26$ cases and $k^{*}=R+1$ in
$13$ (check~(C)). Exhaustive search over short binary sequences gives minimal witnesses:
\texttt{AACAC} has $R=3$, $k^{*}=3$; \texttt{AACACC} has $R=2$, $k^{*}=3$; and homopolymer
runs such as \texttt{AAAA} have $k^{*}<R$, because a repeat that always extends the same
way branches nowhere. Theorem~\\ref{Bio1:thm:repeat} explains the first family and
Theorem~\\ref{Bio1:thm:interleave} the second, but we have no closed form for $k^{*}$ in general.

\paragraph{Open.} (a) \emph{Additivity beyond two repeats.} Theorem~\\ref{Bio1:thm:interleave}
and Proposition~\\ref{Bio1:prop:three} identify the correction at $s=1$ with the sandpile group
of the connector digraph $D_w$ of the circular word of occurrences, and show it is not a
function of the chord diagram's intersection rank. We conjecture
$\K(G_k)\cong\bigoplus_i(\Z/r_i)^{s_i-1}\oplus\K(D_w)$ for $m$ repeats of multiplicities
$r_i$ (with $D_w$ then carrying multiplicities); it is proved for $m\le2$ and verified for
$m=3$, $r_i=2$, $s\le2$. (b) \emph{Higher multiplicity.} We have not
computed the interleaving correction for repeats of multiplicity $r>2$. (c)
\emph{Double-strandedness.} Real assembly graphs are bidirected. This is a different
combinatorial object, and we warn against the tempting decomposition of $\K$ into
$\pm1$-eigenspaces of the reverse-complement involution: an involution splits an abelian
group only after $2$ is inverted, and the $2$-torsion --- which is where all of our
examples live, and where inverted repeats would show up --- is exactly what such a
splitting discards. The tools that survive are the signed and $\Z/2$-voltage Laplacians
\cite{Zaslavsky,RT}, with one caveat: reverse complementation reverses arc direction, so
the passage to the bidirected graph is not a covering of digraphs and the covering theory
of \cite{RT} does not apply as it stands. (d) \emph{Errors and coverage.} We assume an exact $k$-mer spectrum
with true multiplicities.

\section{Verification}\label{Bio1:sec:battery}

The script \texttt{dbg\_sandpile.py} runs $30$ checks in about two minutes, all passing; its
output is archived as \texttt{dbg\_sandpile\_output.txt}. Arithmetic is exact except in the
two rows marked \emph{sampled}.

\begin{table}[ht]
\centering\small
\begin{tabular}{clcl}
\toprule
key & check & \# & mode\\
\midrule
(T) & $\DB(\sigma,k{+}1)=L(\DB(\sigma,k))$; $|\K|=\sigma^{\sigma^{k}-k-1}$ & 8/8 & exact\\
(L) & $\K(\DB(2,n))$ reproduces Levine \cite{Levine}, $n\le5$ & 4/4 & exact\\
(E) & BEST \\eqref{Bio1:eq:best} against brute-force arc-rooted enumeration & 6/6 & exact\\
(R) & Theorem~\\ref{Bio1:thm:repeat}, $54$ parameter settings & 1/1 & exact\\
(I) & Theorem~\\ref{Bio1:thm:interleave}: sequences, compacted graphs, symbolic SNF & 3/3 & exact\\
(M) & Proposition~\\ref{Bio1:prop:three}: all $90$ circular words of three repeats, $s=1,2$ & 1/1 & exact\\
(N) & Theorem~\\ref{Bio1:thm:notdeg} witnesses; degree-class splitting rate & 2/2 & exact / sampled\\
(C) & collapse order: statistics; minimal binary witnesses & 2/2 & sampled / exact\\
(P) & Lemma~\\ref{Bio1:lem:contract}: contraction and self-loops & 2/2 & exact\\
(F) & Table~\\ref{Bio1:tab:phix}: $\varphi$X174 & 1/1 & exact\\
\midrule
 & total & 30/30 & \\
\bottomrule
\end{tabular}
\caption{The verification battery (\texttt{dbg\_sandpile.py}).}
\label{Bio1:tab:battery}
\end{table}

\section*{Provenance and changes in version 2}

Unrefereed preprint. The classical inputs are the BEST theorem \cite{AEdB,SmithTutte},
Knuth's recursion \cite{Knuth} and the matrix--tree theorem for digraphs; the critical
groups of complete de Bruijn graphs \cite{Levine,BidkhoriKishore,CHP} are used only as
checks. Proposition~\\ref{Bio1:prop:decomp} is a reading of \\eqref{Bio1:eq:best}; the use of BEST to
count genome reconstructions is due to \cite{KSP}. What we claim as new is
Theorems~\\ref{Bio1:thm:repeat}, \\ref{Bio1:thm:interleave} and \\ref{Bio1:thm:notdeg} with
Corollary~\\ref{Bio1:cor:insufficient}, and Lemma~\\ref{Bio1:lem:contract} with its use in
Sections~\\ref{Bio1:sec:compute}--\\ref{Bio1:sec:real}.

Version 2 adds Theorem~\\ref{Bio1:thm:interleave} (resolving the two-repeat case of the
additivity question left open in version 1), Proposition~\\ref{Bio1:prop:three} (three repeats,
refuting the chord-rank guess), the real-genome computation of
Section~\\ref{Bio1:sec:real}, the counting convention of Remark~\\ref{Bio1:rem:convention}, the
self-loop clause of Lemma~\\ref{Bio1:lem:contract}, the minimal collapse witnesses, and the
warning about eigenspace decompositions in Open~(c). These changes respond to a detailed
reading of version 1, for which we are grateful; the referee's four corrections were all
right and are all incorporated.
